\subsection{先判断有没有液相，再决定蒸气分压}
饱和蒸气压 $p_s(T)$ 是两相平衡条件，不是所有蒸气在该温度下必有的压强。若蒸气总量太少，没有液体留下，蒸气分压应由实际物质的量和体积计算；只有液相存在且达到平衡时，分压才被 $p_s(T)$ 固定。
混有不凝结气体时，理想混合近似下被固定的也是蒸气\emph{分压}，不是总压强。图\ref{ch17:fig:s31-22-5}的水平共存线针对纯物质；含空气的容器即使正在凝结，总压强也未必保持不变。

\begin{origexamplebox}{湿空气等温压缩的凝结起点}\label{ch17:ex:teaching-condensation}
温度 $T$ 下，体积为 $V_0$ 的密闭容器内有理想干空气和水蒸气，初始无液水。干空气分压为 $p_{a0}$，水蒸气分压为 $\phi p_s(T)$，其中 $0<\phi<1$。缓慢等温压缩至 $V=V_0/r$、$r\ge1$。假定气液平衡及时建立，忽略液水体积及空气在液水中的溶解，求开始凝结时的 $r$、总压强和凝结的水的物质的量。
\end{origexamplebox}
\begin{xesolution}
凝结前两种气体各自的物质的量不变，分压都随压缩倍数增长：$p_a=rp_{a0}$、$p_v=r\phi p_s$。达到饱和的条件为 $r\phi=1$，故临界压缩倍数为 $r_*=1/\phi$。
凝结后，水分子从气相转入液相，不能继续对气相水蒸气单独使用定量气体的玻意耳关系。于是
\[
 p(r)=rp_{a0}+\begin{cases}r\phi p_s,&1\le r\le1/\phi,\\p_s,&r\ge1/\phi.\end{cases}
\]
初始水总物质的量为 $n_{w0}=\phi p_sV_0/(RT)$；凝结后气相水为 $n_v=p_sV_0/(rRT)$，所以
\[
 n_l=\begin{cases}0,&1\le r\le1/\phi,\\
 \displaystyle\frac{p_sV_0}{RT}\left(\phi-\frac1r\right),&r\ge1/\phi.
 \end{cases}
\]
临界点两段总压强连续、凝结量从零开始且非负。气液两相的水相加仍为 $n_{w0}$，而干空气始终满足 $p_aV=p_{a0}V_0$。总压强在凝结后仍增长，增长来自空气分压。若压缩到液体体积不再可忽略，就必须用扣除液体后的气相空间重算，本式不再适用。
\end{xesolution}

\subsection{从共存线斜率到沸点变化}
沸腾需考虑液体内部气泡的机械平衡。在液柱静压和曲率附加压强均可忽略的平面近似下，沸点满足 $p_s(T_b)=p_{\rm ext}$；它并非脱离外压的固定温度。
前面的共存线斜率确定了改变外压后沸点的变化方向：汽化潜热为正、蒸气比体积更大，所以减压使沸点降低。

\begin{origexamplebox}{减压到一半时的沸点}\label{ch17:ex:teaching-boiling}
某纯液体在外压 $p_1$ 下的沸点为 $T_1=360\,\mathrm K$。在所讨论温区内，汽化摩尔潜热近似为 $L_m=40.0\,\mathrm{kJ/mol}$，蒸气服从理想气体关系，液相比体积可忽略。求外压降为 $p_1/2$ 时的沸点，取 $R=8.314\,\mathrm{J/(mol\,K)}$。
\end{origexamplebox}
\begin{xesolution}
将摩尔潜热直接代入共存线积分式，不能再乘一次摩尔质量：
\[
 \ln\frac{p_2}{p_1}=\frac{L_m}{R}\left(\frac1{T_1}-\frac1{T_2}\right).
\]
本题左边为 $-\ln2$，移项得
\[
 \frac1{T_2}=\frac1{360\,\mathrm K}+\frac{R\ln2}{L_m},
 \qquad T_2\simeq342\,\mathrm K.
\]
减压后温度下降，与共存曲线正斜率一致；$R/L_m$ 的单位为 $\mathrm{K^{-1}}$。若压强变化很小，可线性化为 $\Delta T\simeq RT_1^2\ln(p_2/p_1)/L_m$；本题用倒温度形式直接求解，保留了这一模型积分结果的非线性。近临界区、潜热显著变化或液柱很高时，须恢复被忽略的条件。
\end{xesolution}
